% Generated from locale/tr/content/first-order-logic/natural-deduction/proving-things-quant.tex; do not hand-edit.


\olfileid[tr]{fol}{ntd}{prq}

\olsection{Niceleyiciler İçeren \usetoken{P}{derivation}}

\begin{ex}
Niceleyicileri ele alırken özdeğişken koşulunu ihlal etmemeye dikkat
etmeliyiz; bazen bunun için belirli çıkarımları uygulama sıramızı
değiştirmemiz gerekir. Genel olarak özdeğişken koşuluna bağlı kuralları önce
ele almaya çalışmak yararlıdır (tamamlanmış ispatta bu kurallar daha aşağıda
bulunur).

!!a{derivation} nasıl verilir bakalım; hedef !!{formula},
$\lexists[x][\lnot !A(x)] \lif \lnot \lforall[x][!A(x)]$ olsun. Her zamanki
gibi şöyle başlarız:
\begin{prooftree}
\AxiomC{}
\UnaryInfC{$\lexists[x][\lnot !A(x)]\lif \lnot \lforall[x][!A(x)]$}
\end{prooftree}
Son adımı \Intro{\lif} kuralıyla gerekçelendirmek için ne gerektiğini
yazarak başlarız.
\begin{prooftree}
\AxiomC{$\Discharge{\lexists[x][\lnot !A(x)]}{1}$}
\DeduceC{$\lnot \lforall[x][!A(x)]$}
\DischargeRule{\Intro{\lif}}{1}
\UnaryInfC{$\lexists[x][\lnot !A(x)]\lif \lnot \lforall[x][!A(x)]$}
\end{prooftree}
$\lnot \lforall[x][!A(x)]$ ifadesine uygulanacak belirgin bir kural
olmadığından, \Elim{\lexists} kuralını kullanabileceğimiz biçimde
!!{derivation} kurarak ilerleyeceğiz. Burada özdeğişken koşuluna dikkat
etmeli ve $\lexists[x][\lnot !A(x)]$ öncülünde, sonuçta veya ilgili
kapatılmamış varsayımlarda geçmeyen bir sabit seçmeliyiz. (!!{constant}s
kümesi boş olduğundan herhangi bir seçim uygundur.)
\begin{prooftree}
\AxiomC{$\Discharge{\lexists[x][\lnot !A(x)]}{1}$}
\AxiomC{$\Discharge{\lnot !A(a)}{2}$}
\DeduceC{$\lnot \lforall[x][!A(x)]$}
\DischargeRule{\Elim{\lexists}}{2}
\BinaryInfC{$\lnot \lforall[x][!A(x)]$}
\DischargeRule{\Intro{\lif}}{1}
\UnaryInfC{$\lexists[x][\lnot !A(x)] \lif \lnot \lforall[x][!A(x)]$}
\end{prooftree}
$\lnot \lforall[x][!A(x)]$ ifadesini türetmek için \Intro{\lnot} kuralını
kullanmayı deneyeceğiz: bunun için, muhtemelen ek varsayım olarak
$\lforall[x][!A(x)]$'i kullanarak bir çelişki türetmemiz gerekir. Elbette bu
çelişki, \Elim{\lexists} çıkarımı tarafından kapatılacak $\lnot !A(a)$
varsayımını içerebilir. Kurulumu şöyle yapabiliriz:
\begin{prooftree}
\AxiomC{$\Discharge{\lexists[x][\lnot !A(x)]}{1}$}
\AxiomC{$\Discharge{\lnot !A(a)}{2}, \Discharge{\lforall[x][!A(x)]}{3}$}
\DeduceC{$\lfalse$}
\DischargeRule{\Intro{\lnot}}{3}
\UnaryInfC{$\lnot \lforall[x][!A(x)]$}
\DischargeRule{\Elim{\lexists}}{2}
\BinaryInfC{$\lnot \lforall[x][!A(x)]$}
\DischargeRule{\Intro{\lif}}{1}
\UnaryInfC{$\lexists[x][\lnot !A(x)]\lif \lnot \lforall[x][!A(x)]$}
\end{prooftree}
Bir çelişkiye ulaşmaya yaklaşmış görünüyoruz. Uygulanacak en kolay kural,
özdeğişken koşulu bulunmayan \Elim{\lforall} kuralıdır. Evrensel
niceleyicinin bağladığı $x$ yerine istediğimiz herhangi bir terimi
koyabildiğimiz için, çelişkiye ulaşmak amacıyla $a$'yı kullanmaya devam
etmek en uygunudur.
\begin{prooftree}
\AxiomC{$\Discharge{\lexists[x][\lnot !A(x)]}{1}$}
\AxiomC{$\Discharge{\lnot !A(a)}{2}$}
\AxiomC{$\Discharge{\lforall[x][!A(x)]}{3}$}
\RightLabel{\Elim{\lforall}}
\UnaryInfC{$!A(a)$}
\RightLabel{\Elim{\lnot}}
\BinaryInfC{$\lfalse$}
\DischargeRule{\Intro{\lnot}}{3}
\UnaryInfC{$\lnot \lforall[x][!A(x)]$}
\DischargeRule{\Elim{\lexists}}{2}
\BinaryInfC{$\lnot \lforall[x][!A(x)]$}
\DischargeRule{\Intro{\lif}}{1}
\UnaryInfC{$\lexists[x][\lnot !A(x)]\lif \lnot \lforall[x][!A(x)]$}
\end{prooftree}

Özellikle niceleyicilerle uğraşırken, bu aşamada özdeğişken koşulunun ihlal
edilmediğini yeniden denetlemek önemlidir. Uyguladığımız ve özdeğişken
koşuluna bağlı olan tek kural \Elim{\lexists} olduğundan ve özdeğişken~$a$
çıkarımın bağlı olduğu kapatılmamış varsayımların hiçbirinde geçmediğinden bu, doğru bir
!!{derivation} olur.
\end{ex}

\begin{ex}
Bazen !!a{formula}, başka !!{formula}s kullanılarak türetilebilir. Bu durumlarda
kapatılmamış varsayımlarımız olabilir. Hedefimizin yanı sıra varsayımlarımızı
da izlemek önemlidir.

Şimdi !!a{derivation} kuracağız; hedef !!{formula} $\lexists[x][!C(x,b)]$,
varsayımlar ise $\lexists[x][(!A(x) \land !B(x))]$ ve
$\lforall[x][(!B(x) \lif !C(x,b))]$ olsun. Her zamanki gibi sonucu alta
yazarız.
\begin{prooftree}
\AxiomC{}
\UnaryInfC{$\lexists[x][!C(x,b)]$}
\end{prooftree}

Kullanabileceğimiz iki öncül vardır. İlkini kullanmak, yani
$\lexists[x][(!A(x) \land !B(x))]$ ifadesinden $\lexists[x][!C(x, b)]$ için
!!a{derivation} bulmaya çalışmak üzere \Elim{\lexists} kuralını kullanırız.
Bu kuralın özdeğişken koşulu olduğundan önce onu uygularız. Şunu elde ederiz:
\begin{prooftree}
\AxiomC{$\lexists[x][(!A(x) \land !B(x))]$}
\AxiomC{$\Discharge{!A(a) \land !B(a)}{1}$}
\DeduceC{$\lexists[x][!C(x,b)]$}
\DischargeRule{\Elim{\lexists}}{1}
\BinaryInfC{$\lexists[x][!C(x,b)]$}
\end{prooftree}
Üzerinde çalıştığımız iki varsayımın ortak bileşeni $!B$'dir. Bu noktada
\Elim{\land} uygulayıp $!B(a)$'yı ayırmak yararlı olabilir.
\begin{prooftree}
\AxiomC{$\lexists[x][(!A(x) \land !B(x)])$}
\AxiomC{$\Discharge{!A(a) 
\land !B(a)}{1}$}
\RightLabel{\Elim{\land}}
\UnaryInfC{$!B(a)$}
\DeduceC{$\lexists[x][!C(x,b)]$}
\DischargeRule{\Elim{\lexists}}{1}
\BinaryInfC{$\lexists[x][!C(x,b)]$}
\end{prooftree}

Kullanabileceğimiz ikinci varsayım $\lforall[x][(!B(x) \lif !C(x,b))]$'dir.
Özdeğişken koşulu bulunmadığından \Elim{\lforall} kullanıp $x$'i
!!{constant}~$a$ ile örnekleyerek $!B(a) \lif !C(a, b)$ elde edebiliriz.
Artık hem $!B(a) \lif !C(a,b)$ hem de $!B(a)$ elimizdedir. Sonraki adımımız
\Elim{\lif} kuralının doğrudan bir uygulaması olmalıdır.
\begin{prooftree}
\AxiomC{$\lexists[x][(!A(x) \land !B(x))]$}
\AxiomC{$\lforall[x][(!B(x) \lif !C(x,b))]$}
\RightLabel{\Elim{\lforall}}
\UnaryInfC{$!B(a) \lif !C(a,b)$}
\AxiomC{$\Discharge{!A(a) 
\land !B(a)}{1}$}
\RightLabel{\Elim{\land}}
\UnaryInfC{$!B(a)$}
\RightLabel{\Elim{\lif}}
\BinaryInfC{$!C(a,b)$}
\DeduceC{$\lexists[x][!C(x,b)]$}
\DischargeRule{\Elim{\lexists}}{1}
\insertBetweenHyps{\hspace{-5em}}
\BinaryInfC{$\lexists[x][!C(x,b)]$}
\end{prooftree}
Hedefe çok yaklaştık!{} Bir \Intro{\lexists} uygulamasıyla hedefimize
ulaşırız.
\begin{prooftree}
\AxiomC{$\lexists[x][(!A(x) \land !B(x))]$}
\AxiomC{$\lforall[x][(!B(x) \lif !C(x,b))]$}
\RightLabel{\Elim{\lforall}}
\UnaryInfC{$!B(a) \lif !C(a,b)$}
\AxiomC{$\Discharge{!A(a) 
\land !B(a)}{1}$}
\RightLabel{\Elim{\land}}
\UnaryInfC{$!B(a)$}
\RightLabel{\Elim{\lif}}
\BinaryInfC{$!C(a,b)$}
\RightLabel{\Intro{\lexists}}
\UnaryInfC{$\lexists[x][!C(x,b)]$}
\DischargeRule{\Elim{\lexists}}{1}
\insertBetweenHyps{\hspace{-5em}}
\BinaryInfC{$\lexists[x][!C(x,b)]$}
\end{prooftree}
Her adımda özdeğişken koşullarının ihlal edilmemesini sağladığımız için
bunun doğru bir !!{derivation} olduğundan emin olabiliriz.
\end{ex}

\begin{ex}
$\lforall[x][!A(x)] \lif \lexists[y][!B(y)]$ ve
$\lnot\lexists[y][!B(y)]$ varsayımlarından bir !!a{derivation} verin; hedef
!!{formula} $\lnot\lforall[x][!A(x)]$ olsun. Her zamanki gibi hedef
!!{formula} ifadesini alta yazarak başlarız.
\begin{prooftree}
\AxiomC{}
\UnaryInfC{$\lnot\lforall[x][!A(x)]$}
\end{prooftree}
Bu !!{derivation} içindeki son satır bir değilleme olduğundan \Intro{\lnot}
kullanmayı deneyelim. Bunun için bir çelişkiyi !!{derive} üzere ne
yapabileceğimizi bulmamız gerekir.
\begin{prooftree}
\AxiomC{$\Discharge{\lforall[x][!A(x)]}{1}$}
\DeduceC{$\lfalse$}
\DischargeRule{\Intro{\lnot}}{1}
\UnaryInfC{$\lnot\lforall[x][!A(x)]$}
\end{prooftree}
Buraya kadar iyi. \Elim{\lforall} kullanabiliriz; ama bunun hedefimize
ulaşmaya yardımcı olacağı açık değildir. Onun yerine varsayımlarımızdan
birini kullanalım. $\lforall[x][!A(x)] \lif \lexists[y][!B(y)]$, yanında
$\lforall[x][!A(x)]$ ile birlikte \Elim{\lif} kuralını kullanmamızı sağlar.
\begin{prooftree}
\AxiomC{$\lforall[x][!A(x)] \lif \lexists[y][!B(y)]$}
\AxiomC{$\Discharge{\lforall[x][!A(x)]}{1}$}
\RightLabel{\Elim{\lif}}
\BinaryInfC{$\lexists[y][!B(y)]$}
\DeduceC{$\lfalse$}
\DischargeRule{\Intro{\lnot}}{1}
\UnaryInfC{$\lnot\lforall[x][!A(x)]$}
\end{prooftree}
Artık kullanabileceğimiz son bir varsayım var; bu varsayım \Elim{\lnot}
kullanarak çelişkiye ulaşmamızı sağlayacak gibi görünüyor.
\begin{prooftree}
\AxiomC{$\lnot\lexists[y][!B(y)]$}
\AxiomC{$\lforall[x][!A(x)] \lif \lexists[y][!B(y)]$}
\AxiomC{$\Discharge{\lforall[x][!A(x)]}{1}$}
\RightLabel{\Elim{\lif}}
\BinaryInfC{$\lexists[y][!B(y)]$}
\RightLabel{\Elim{\lnot}}
\BinaryInfC{$\lfalse$}
\DischargeRule{\Intro{\lnot}}{1}
\UnaryInfC{$\lnot\lforall[x][!A(x)]$}
\end{prooftree}
\end{ex}

\begin{prob}
Aşağıdakileri gösteren !!{derivation}s verin:
\begin{enumerate}
\item $\Proves (\lforall[x][!A(x)] \land \lforall[y][!B(y)]) \lif
\lforall[z][(!A(z) \land !B(z))]$.
\item $\Proves (\lexists[x][!A(x)] \lor \lexists[y][!B(y)]) \lif
\lexists[z][(!A(z) \lor !B(z))]$.
\item $\lforall[x][(!A(x) \lif !B)] \Proves \lexists[y][!A(y)] \lif !B$.
\item $\lforall[x][\lnot !A(x)] \Proves \lnot\lexists[x][!A(x)]$.
\item $\Proves \lnot\lexists[x][!A(x)] \lif \lforall[x][\lnot !A(x)]$.
\item $\Proves \lnot\lexists[x][\lforall[y][((!A(x,y) \lif \lnot
!A(y,y)) \land (\lnot !A(y,y) \lif !A(x,y)))]]$.
\end{enumerate}
\end{prob}

\begin{prob}
Aşağıdakileri gösteren !!{derivation}s verin:
\begin{enumerate}
\item $\Proves \lnot\lforall[x][!A(x)] \lif \lexists[x][\lnot!A(x)]$.
\item $(\lforall[x][!A(x)] \lif !B) \Proves \lexists[y][(!A(y) \lif !B)]$.
\item $\Proves \lexists[x][(!A(x) \lif \lforall[y][!A(y)])]$.
\end{enumerate}
(Bunların hepsi $\FalseCl$ kuralını gerektirir.)
\end{prob}

